\documentclass[conference,a4paper]{IEEEtran}
\usepackage[utf8]{inputenc}
\usepackage[T1]{fontenc}
\usepackage[brazil]{babel}
\usepackage{amsmath,amssymb,graphicx,booktabs}
\usepackage[hidelinks]{hyperref}
\setlength{\textfloatsep}{8pt plus 2pt minus 2pt}
\setlength{\floatsep}{7pt plus 2pt minus 2pt}
\newcommand{\U}{\mathrm{m}/\mathrm{s}^{2}}
\newcommand{\PU}{(\mathrm{m}/\mathrm{s}^{2})^{2}}
\newcommand{\fig}[3]{\begin{figure}[!ht]\centering\includegraphics[width=\columnwidth]{figuras/#1.png}\caption{#2}\label{#3}\end{figure}}
\renewcommand{\IEEEkeywordsname}{Palavras-chave}
\begin{document}
\title{Análise temporal e espectral de vibrações de drones\\[0.3em]
\large Primeira entrega de Sinais e Sistemas}
\author{\IEEEauthorblockN{Ana Maria Cunha \quad Arthur Britto \quad Mateus Ferreira Borges}
\IEEEauthorblockA{CIN0020 --- Sinais e Sistemas, 2026.2}}
\maketitle
\begin{abstract}
Este trabalho tem como objetivo analisar padrões de vibração de drones a partir de sinais de aceleração do dataset DronePropB e comparar condições de hélices saudáveis e danificadas nos domínios do tempo e da frequência. O dataset contém 111 experimentos em arquivos MATLAB, todos com uma série de um sensor piezoelétrico e 108 também com três séries de uma IMU (Unidade de Medição Inercial). A partir de registros com configurações operacionais comparáveis, são extraídas características temporais e espectrais para investigar a intensidade e a distribuição em frequência das vibrações. O estudo inclui os efeitos da duração dos trechos e da taxa de amostragem na análise espectral, além de uma investigação das diferenças entre os três eixos da IMU.

\end{abstract}
\begin{IEEEkeywords}
vibração, drones, FFT, amostragem, características de sinais.
\end{IEEEkeywords}

\section{Do problema aos sinais selecionados}
Danos nas hélices podem alterar tanto a intensidade quanto a distribuição em frequência das vibrações de um drone. Nosso objetivo é verificar o que cada domínio revela sobre essas mudanças, sem pressupor que toda falha produza vibração mais intensa. O percurso acompanha o \href{https://colab.research.google.com/drive/1FdcHGuE1c6JbiSro7K0QYuGb-4w69KOJ?usp=sharing}{\textbf{Colab do projeto}}: conhecer os arquivos, selecionar condições comparáveis, extrair medidas, interpretar os espectros e investigar os eixos da IMU. Esta entrega cobre as análises solicitadas; o treinamento de classificadores fica para a segunda etapa.

\subsection{Dataset e recorte de comparação}
O DronePropB reúne ensaios em solo com hélices saudáveis e defeituosas~\cite{dataset}. O inventário do notebook encontrou 111 arquivos MATLAB: todos têm o piezo \texttt{AccData.Vib\_Data}; 108 têm também a IMU triaxial \texttt{IMUData.Vib\_Data} e a tensão da bateria, usada apenas na verificação da Seção~\ref{sec:comp}. As taxas declaradas são 10~kHz no piezo e 1~kHz na IMU. Adotamos aceleração em $\U$, conforme orientação registrada no notebook; a média próxima de $g$ no eixo $z$ da IMU é consistente com essa unidade.

Os nomes codificam condição/grupo de falha (F), severidade (SV), configuração de velocidade (SP), código C e, nos saudáveis, repetição (R). F0/SV0 representa a condição saudável; F1--F3 e SV1--SV3 identificam grupos de falha e severidades, sem atribuir percentuais de dano. O significado físico de C, a rotação em RPM e a orientação dos sensores não foram estabelecidos. Portanto, mantemos SP e C constantes e não atribuímos causas mecânicas específicas aos picos.

A comparação usa \textbf{SP3/C1}, com os 12 registros da Tabela~\ref{tab:registros}. Todos os piezos têm 600\,000 amostras e 60~s. A escolha mantém sensor, taxa, duração e códigos operacionais constantes, cobrindo todas as falhas e severidades. O hash do piezo revelou nove pares de sinais saudáveis idênticos entre C1 e C2 no inventário geral (102 sinais distintos em 111 arquivos); usar somente C1 evita contar essas cópias como novos experimentos.

\begin{table}[!ht]
\caption{Seleção SP3/C1. IDs sem extensão .mat. O piezo é uniaxial, a 10~kHz, com 600\,000 amostras por registro. A IMU tem três eixos a 1~kHz; $N_i$ é por eixo e sua duração nominal é $N_i/1000$ segundos.}
\label{tab:registros}\centering\footnotesize\setlength{\tabcolsep}{4pt}
\begin{tabular}{@{}llr@{}}\toprule
Registro & Rótulo & $N_i$\\\midrule
\texttt{F0\_SV0\_SP3\_C1\_R1} & Saudável & 81\,709\\
\texttt{F0\_SV0\_SP3\_C1\_R2} & Saudável & 84\,400\\
\texttt{F0\_SV0\_SP3\_C1\_R3} & Saudável & 81\,933\\
\texttt{F1\_SV1\_SP3\_C1} & Danificada & 81\,639\\
\texttt{F1\_SV2\_SP3\_C1} & Danificada & 80\,828\\
\texttt{F1\_SV3\_SP3\_C1} & Danificada & 80\,784\\
\texttt{F2\_SV1\_SP3\_C1} & Danificada & 80\,613\\
\texttt{F2\_SV2\_SP3\_C1} & Danificada & 81\,196\\
\texttt{F2\_SV3\_SP3\_C1} & Danificada & 80\,706\\
\texttt{F3\_SV1\_SP3\_C1} & Danificada & 82\,245\\
\texttt{F3\_SV2\_SP3\_C1} & Danificada & 82\,001\\
\texttt{F3\_SV3\_SP3\_C1} & Danificada & 81\,033\\\bottomrule
\end{tabular}
\end{table}

\subsection{Verificação, preparação e janelas}
A inspeção dos registros completos e da média móvel de 0,5~s não motivou tratamento do piezo. Usamos os \textbf{dados brutos, sem filtro}, preservando o procedimento atual do Colab. A conferência local confirmou valores finitos, comprimentos compatíveis entre sinal e tempo e timestamps crescentes e regulares a 0,1~ms nos 12 piezos. Isso verifica a base temporal, mas não certifica calibração ou ausência de saturação.

Cada registro é dividido em 60 janelas de 1~s (10\,000 amostras), sem sobreposição nem sobras: \textbf{180 saudáveis e 540 danificadas}. Cada linha de características conserva arquivo, condição, F, SV e índice da janela no notebook. A FFT usa esses mesmos trechos. As 720 janelas vêm de apenas 12 aquisições e não são experimentos independentes; os danificados também reúnem grupos e severidades diferentes. Manter SP/C constantes melhora a comparação, mas não demonstra que o dano seja a única causa das diferenças.

\clearpage
\section{O que a intensidade e os extremos mostram}
A Fig.~\ref{fig:tempo} compara o mesmo sensor em F0/SV0/R1 e F1/SV3, ambos SP3/C1. O registro completo mostra variações de amplitude; o zoom de 50~ms evidencia oscilações rápidas e a maior amplitude desse exemplo danificado. Os dois sinais compartilham as escalas em cada painel. Essa comparação ilustra um par de registros; a conclusão sobre o conjunto depende das distribuições de todas as janelas.

\fig{01_notebook_tempo}{Piezo bruto: F0/SV0/R1 (saudável) e F1/SV3 (danificada), SP3/C1. Acima, 60~s; abaixo, os primeiros 50~ms. Imagem original do notebook; o título ainda diz ``filtrado, mas o código usa os sinais brutos, sem filtro.}{fig:tempo}

\subsection{Definições e convenções}
Para uma janela $x[n]$, com $N$ amostras, usamos as definições da Tabela~\ref{tab:def-tempo}. O desvio padrão é populacional, com divisor $N$. Média representa nível contínuo; RMS e potência descrevem intensidade global; desvio mede flutuação em torno da média; os picos descrevem extremos. O fator de crista mede quanto o maior extremo se destaca do RMS. Potência do sinal tem unidade quadrática de aceleração, e não watts.

\begin{table}[!ht]
\caption{Características temporais por janela. $A$ denota $\U$; $A^2$ denota $\PU$.}
\label{tab:def-tempo}\centering\small\setlength{\tabcolsep}{4pt}
\begin{tabular}{@{}lll@{}}\toprule
Medida & Definição & Unidade\\\midrule
Média & $\mu=N^{-1}\sum_n x[n]$ & $A$\\
Potência & $P=N^{-1}\sum_n x[n]^2$ & $A^2$\\
RMS & $\sqrt{P}$ & $A$\\
Desvio padrão & $\sigma=\sqrt{N^{-1}\sum_n(x[n]-\mu)^2}$ & $A$\\
Pico absoluto & $a_p=\max_n |x[n]|$ & $A$\\
Pico a pico & $\max_n x[n]-\min_n x[n]$ & $A$\\
Fator de crista & $a_p/\mathrm{RMS}$ & 1\\\bottomrule
\end{tabular}
\end{table}

\begin{table}[!ht]
\caption{Características temporais: média $\pm$ desvio entre janelas. Unidades conforme Tabela~\ref{tab:def-tempo}.}
\label{tab:tempo}\centering\small\setlength{\tabcolsep}{4pt}
\begin{tabular}{@{}lrr@{}}\toprule
Medida & Saudável & Danificada\\\midrule
Média & $-0,00144\pm0,00125$ & $-0,00170\pm0,00164$\\
RMS & $0,2729\pm0,0495$ & $0,3046\pm0,1186$\\
Potência & $0,07691\pm0,02640$ & $0,10682\pm0,07998$\\
Desvio padrão & $0,2729\pm0,0495$ & $0,3046\pm0,1186$\\
Pico absoluto & $1,008\pm0,204$ & $1,107\pm0,356$\\
Pico a pico & $1,886\pm0,380$ & $2,109\pm0,685$\\
Fator de crista & $3,708\pm0,443$ & $3,755\pm0,509$\\\bottomrule
\end{tabular}
\end{table}

\fig{02_notebook_tempo_box}{Distribuições das medidas temporais nas 720 janelas. Caixas: intervalo interquartil e mediana; bigodes: até 1,5 vezes esse intervalo. Os pontos além dos bigodes também são apresentados; nenhuma janela foi excluída desses cálculos.}{fig:tempo-box}

\subsection{Mais vibração nem sempre significa falha}
A Tabela~\ref{tab:tempo} mostra RMS médio cerca de 12\% maior nos danificados, mas dispersão 2,4 vezes maior. Os boxplots (Fig.~\ref{fig:tempo-box}) têm forte sobreposição; a mediana danificada é até menor. No detalhamento por grupo do Colab, as medianas do RMS são aproximadamente 0,46~$\U$ em F1, 0,25 em F2 e 0,19 em F3, contra 0,28 no saudável. Portanto, a média agrupada esconde comportamentos distintos. O fator de crista quase não muda entre as duas classes, embora varie entre grupos de falha.

A identidade $P=\mathrm{RMS}^2=\sigma^2+\mu^2$ explica por que RMS, potência e desvio trazem informação semelhante no piezo, cuja média é pequena. O exemplo central de 1~s da IMU $z$ no Colab evidencia o efeito oposto: com $\mu=9,69$ e $\sigma=7,35$~$\U$, o RMS é 12,17~$\U$; removida a média, cai a 7,35~$\U$. A gravidade eleva o RMS sem aumentar a flutuação. Os trechos centrais dos dois sensores não são considerados sincronizados.

\clearpage
\section{Onde a potência está no espectro}
\subsection{FFT, magnitude e potência}
Para cada uma das mesmas 720 janelas, retiramos a média e aplicamos Hann antes da FFT real. O eixo é $f_k=kf_s/N$, de 0 a 5\,000~Hz, com espaçamento $\Delta f=1$~Hz. A magnitude unilateral é $|X_w[k]|/\sum_n w[n]$, dobrada nas frequências com par negativo; DC e Nyquist não são dobradas. A unidade permanece $\U$. A normalização compensa o ganho da Hann para uma senoide centrada em um bin; o Colab a confere com amplitudes 3 e 1 em 120 e 350~Hz. Fora dos bins, o pico pode ser atenuado e espalhado. A Hann reduz vazamento, mas alarga picos: espaçamento da grade não equivale, sozinho, à capacidade de separar componentes.

Para potência, o notebook usa outra FFT: \textbf{sem Hann}, após retirar a média. Define $q_k=|X[k]|^2/N^2$, com a mesma compensação unilateral. Por Parseval, $\sum_k q_k=\sigma^2$; essa potência exclui a componente média. Não se soma o quadrado da magnitude corrigida pela Hann para estimar potência.

\fig{03_notebook_espectros}{Média das magnitudes por condição nas mesmas janelas da análise temporal. Escalas linear (acima) e logarítmica (abaixo), na faixa ampliada exibida pelo notebook. A média reúne três registros saudáveis ou nove danificados.}{fig:espectros}

\subsection{Características e escolha das bandas}
A partir de 5~Hz, a frequência dominante é o bin de maior magnitude e a magnitude dominante é sua altura. Para centroide e entropia, usamos $p_k=M_k^2/\sum_j M_j^2$, sendo $M_k$ a magnitude com Hann. O centroide $\sum_k f_kp_k$ indica onde se concentra o espectro; a entropia $-\sum_k p_k\log_2p_k/\log_2K$ mede sua dispersão entre $K$ bins, de zero a um.

A potência em banda é a soma dos $q_k$ entre limites inclusivo à esquerda e exclusivo à direita. As bandas finais são B1: 5--200~Hz, B2: 200--335~Hz, B3: 335--600~Hz e B4: 600--5\,000~Hz. Elas separam as regiões observadas próximas de 130, 270 e 400~Hz e o conteúdo restante, usando as mesmas fronteiras nas duas condições. Substituem as três bandas preliminares exploradas no início da seção 3.3 do Colab. Frequências abaixo de 5~Hz e o bin exato de Nyquist não entram nas bandas, que por isso não somam exatamente a potência total.

\begin{table}[!ht]
\caption{Características espectrais: média $\pm$ desvio entre janelas. Potência total e bandas em $\PU$; magnitude em $\U$; frequências em Hz; entropia adimensional.}
\label{tab:freq}\centering\small\setlength{\tabcolsep}{3pt}
\begin{tabular}{@{}lrr@{}}\toprule
Medida & Saudável & Danificada\\\midrule
Freq. dominante & $243,64\pm130,40$ & $252,40\pm126,28$\\
Magnitude do pico & $0,2228\pm0,0626$ & $0,2478\pm0,1289$\\
Potência total AC & $0,07690\pm0,02640$ & $0,10681\pm0,07998$\\
Centroide & $357,66\pm51,91$ & $367,99\pm62,28$\\
Entropia & $0,4346\pm0,0662$ & $0,4385\pm0,0542$\\
B1 & $0,03476\pm0,02073$ & $0,04694\pm0,05174$\\
B2 & $0,01037\pm0,00639$ & $0,00781\pm0,00459$\\
B3 & $0,02809\pm0,01264$ & $0,04771\pm0,03098$\\
B4 & $0,00367\pm0,00108$ & $0,00433\pm0,00266$\\\bottomrule
\end{tabular}
\end{table}

\subsection{Picos e redistribuição}
Na Fig.~\ref{fig:espectros}, o saudável tem picos em 133~Hz (0,155~$\U$), 266~Hz (0,039) e 400~Hz (0,073). No danificado, destacam-se 128~Hz (0,097) e 375/382~Hz (cerca de 0,060). Os picos do espectro médio agrupado ficam mais baixos quando os registros têm máximos em bins diferentes; isso não implica menor potência. Relações próximas de dobro/triplo e espaçamentos de aproximadamente 60~Hz são observáveis, mas não identificam, por si só, rotação ou origem física.

A Tabela~\ref{tab:freq} revela aumento de cerca de 70\% em B3 e redução de 25\% em B2. A média das frações de potência por janela passa de 38,0\% para 50,7\% em B3; B1 cai de 43,3\% para 36,3\%, B2 de 13,8\% para 8,5\%, e B4 fica próxima de 5\%. Assim, a diferença inclui uma mudança de distribuição, além da intensidade global. O deslocamento de picos é descritivo deste recorte: mesmo código SP não comprova a mesma RPM efetiva, e a Seção~\ref{sec:comp} mostra que ele acompanha a tensão da bateria.

\clearpage
\section{O que muda ao observar e amostrar o sinal}
\subsection{Duração do trecho e separação dos picos}
A Fig.~\ref{fig:resolucao} usa o início do \textbf{mesmo registro saudável F0/SV0/SP3/C1/R1}, com Hann e a normalização anterior, em duas durações. Para 0,1~s, $N=1\,000$ e $\Delta f=10$~Hz. Para 2~s, $N=20\,000$ e $\Delta f=0,5$~Hz. A mudança vem do comprimento observado, sem inserir amostras artificiais.

\fig{04_notebook_resolucao}{FFT dos primeiros 0,1 e 2~s do mesmo piezo. Acima, trecho de 0,1~s; abaixo, trecho de 2~s, ambos até 600~Hz e com escalas comuns. Os pontos indicam os bins calculados.}{fig:resolucao}

\fig{05_notebook_resolucao_zoom}{Detalhe da região dominante para os mesmos dois trechos, conforme o zoom original do notebook.}{fig:resolucao-zoom}

No trecho curto, a região dominante se espalha entre 130 e 140~Hz, com altura em torno de 0,185~$\U$. No trecho de 2~s, o máximo fica próximo de 134~Hz, em torno de 0,245~$\U$, e aparece estrutura vizinha perto de 136~Hz. A grade curta não permite inspecionar esse detalhe. Frequência fora do bin, vazamento e largura da Hann afetam a posição e a altura observadas; os trechos também não têm exatamente o mesmo conteúdo temporal. Portanto, a diferença não deve ser interpretada como mudança física causada pela FFT.

Trechos longos ajudam a distinguir componentes próximas, mas resumem um intervalo maior e podem misturar variações no tempo. A janela principal de 1~s, com grade de 1~Hz, oferece um compromisso simples entre detalhe espectral e acompanhamento local, sem garantir estacionariedade ou separação de todo par de picos.

\subsection{Redução de taxa e aliasing}
A segunda experiência conserva os primeiros 2~s desse registro e reduz a taxa de 10\,000 para \textbf{200~Hz}, por um fator 50. Primeiro selecionamos uma amostra a cada 50, \textbf{sem filtro}; depois repetimos a redução com \texttt{resample\_poly}, que inclui antialias. A Fig.~\ref{fig:aliasing} compara os dois casos com o original.

\fig{06_notebook_aliasing}{Mesmo trecho de 2~s: original, redução sem filtro e redução com antialias. A linha tracejada marca a nova Nyquist de 100~Hz; os sinais reduzidos só possuem espectro até esse limite.}{fig:aliasing}

Com $f_s'=200$~Hz, a nova Nyquist é 100~Hz. A componente dominante próxima de 134~Hz ultrapassa esse limite e, sem filtro, reaparece aproximadamente em
\[
f_a=|134-200|=66\ \mathrm{Hz}.
\]
Outras componentes acima de 100~Hz também se rebatem para 0--100~Hz e se misturam ao conteúdo verdadeiro. Logo, simplesmente reduzir o número de amostras pode produzir uma interpretação errada das frequências presentes.

O filtro antialias atenua componentes altas \emph{antes} da redução. O espectro reduzido filtrado preserva aproximadamente o conteúdo abaixo do corte, com transição e atenuação próprias do filtro, e suprime o falso pico forte em 66~Hz. A componente de 134~Hz é perdida: o filtro evita o rebatimento, mas não torna possível representar essa frequência a 200~Hz. A taxa deve ser escolhida a partir da faixa que se deseja preservar, com margem para a transição do filtro; a Nyquist original de 5~kHz dá uma faixa muito maior para esta análise.

\clearpage
\section{Comparação entre os domínios}
% (hoje sem conteúdo) no Overleaf.
% Requer figuras/09_notebook_comparacao.png (está nesta mesma pasta do repasse).
% Usa apenas o que já existe no preâmbulo: booktabs e as macros \U e \fig.

\section{Comparação entre os domínios}\label{sec:comp}
As características dos dois domínios vêm das mesmas 720 janelas. Para medir quanto cada uma separa as condições, usamos o $d$ de Cohen (diferença entre as médias, em desvios padrão) e a AUC (probabilidade de uma janela danificada ter valor maior que uma saudável; 0,5 indica nenhuma separação). Acrescentamos a \textbf{fração} da potência em cada banda (potência da banda dividida pela potência total da janela), que descreve a forma do espectro sem depender da intensidade.

\begin{table}[!ht]
\caption{Separação entre condições: $d$ de Cohen e AUC com as nove falhas reunidas, e $d$ de cada grupo de falha em relação às janelas saudáveis.}
\label{tab:comp}\centering\small\setlength{\tabcolsep}{3.5pt}
\begin{tabular}{@{}lrrrrr@{}}\toprule
 & \multicolumn{2}{c}{Reunidas} & \multicolumn{3}{c}{$d$ por grupo}\\\cmidrule(lr){2-3}\cmidrule(l){4-6}
Característica & $d$ & AUC & F1 & F2 & F3\\\midrule
RMS & $+0,30$ & 0,51 & $+1,34$ & $-0,16$ & $-0,27$\\
Fator de crista & $+0,09$ & 0,52 & $-0,88$ & $+0,41$ & $+0,61$\\\midrule
Freq. dominante & $+0,07$ & 0,31 & $-1,25$ & $+1,07$ & $+0,23$\\
Centroide & $+0,17$ & 0,56 & $-1,61$ & $+1,57$ & $+0,71$\\
Entropia & $+0,07$ & 0,53 & $-0,92$ & $+0,56$ & $+0,49$\\
Fração em B2 & $-1,38$ & 0,27 & $-1,45$ & $-0,79$ & $-0,98$\\
Fração em B3 & $+0,69$ & 0,69 & $-0,85$ & $+3,06$ & $+1,30$\\\bottomrule
\end{tabular}
\end{table}

Com as nove falhas reunidas, o RMS não separa as condições ($d=0,30$; AUC $=0,51$). A Tabela~\ref{tab:comp} mostra o motivo: os grupos de falha têm efeitos de sentidos opostos. F1 aumenta a intensidade (RMS: $d=+1,34$) e concentra a potência nas baixas frequências; F2 mantém a intensidade e tem muito mais potência em B3 ($d=+3,06$); F3 fica entre os dois. Nenhum domínio é superior em todos os casos: as medidas de intensidade destacam F1, e as de distribuição da potência destacam F2 e F3.

\fig{09_notebook_comparacao}{(a) RMS (tempo) contra a fração da potência em B3 (frequência); cada ponto é uma janela de 1~s. (b) Frequência do pico entre 100 e 160~Hz contra a tensão média da bateria; cada ponto é um registro. Os quadrados são os nove registros danificados de SP3/C2, usados apenas nesta verificação.}{fig:comp}

\subsection{Mesma intensidade, espectros diferentes}
Na Fig.~\ref{fig:comp}(a), janelas com o mesmo RMS têm distribuições de potência muito distintas. Entre as janelas com RMS na faixa central das saudáveis (0,24 a 0,31~$\U$), a fração média em B3 é 36\% nas saudáveis, 17\% em F1 e 74\% em F2. O tempo mede quanto o sinal vibra; a frequência mostra onde essa potência está.

\subsection{A posição dos picos acompanha a bateria}
O deslocamento do pico principal (133~Hz nos saudáveis; 128~Hz, em média, nos danificados) é invisível no tempo: equivale a 0,3~ms no período. Porém, a frequência desse pico acompanha a tensão da bateria registrada nos arquivos (Fig.~\ref{fig:comp}(b)): $r=0,992$ nos 12 registros e $r=0,995$ ao incluir os nove registros danificados de SP3/C2, cerca de 7~Hz/V. Os saudáveis foram gravados com 15,4 a 15,5~V e os danificados com 14,3 a 15,3~V. A mesma falha F1/SV1 tem o pico em 130,1~Hz a 14,97~V (C1) e em 135,1~Hz a 15,65~V (C2), acima dos saudáveis. Com estes dados, o deslocamento é explicado pela tensão da bateria e não pode ser atribuído ao defeito. O padrão é compatível com variação da rotação dos motores, que o dataset não informa.

\subsection{Redundância e características para a classificação}
RMS, desvio padrão, potência média, potência espectral total, picos e magnitude do pico dominante medem a mesma intensidade ($r\geq0,92$ entre janelas), e as potências absolutas por banda a acompanham ($r$ de 0,76 a 0,99 com o RMS). A fração em B1 e o centroide são redundantes com a fração em B3 ($r=-0,97$ e $+0,90$). Para a segunda entrega, propomos: modelo temporal com \textbf{RMS e fator de crista}; modelo espectral com \textbf{frações em B2 e B3 e entropia}; modelo combinado com as cinco. A frequência dominante fica fora por alternar entre dois picos e acompanhar a bateria. A fração em B3 também se relaciona com a tensão, de forma mais fraca ($r\approx-0,6$ por registro): registramos isso como risco de confusão a acompanhar na classificação.

\section{Investigação adicional: a direção faz diferença?}
\textbf{Pergunta:} as vibrações nos eixos $x$, $y$ e $z$ da IMU têm as mesmas características, e a diferença saudável--danificada aparece igualmente nas três direções? A pergunta verifica o que se perde ao observar uma única direção e qual eixo merece atenção em uma aplicação futura.

\subsection{Método e cuidados}
Usamos as IMUs dos mesmos 12 registros SP3/C1, a 1~kHz, com Nyquist de 500~Hz. Diferentemente do piezo, aplicamos passa-altas Butterworth de ordem 4, corte em 5~Hz, por \texttt{filtfilt} no registro completo de cada eixo, antes do janelamento. Isso atenua gravidade e variações lentas, sem deslocamento de fase; a aplicação nos dois sentidos também modifica a resposta de magnitude. O filtro não retira apenas gravidade e pode afetar bordas.

As janelas têm 1~s, ou 1\,000 amostras. Para reduzir a influência dos trechos iniciais de baixa atividade, o Colab retém, por registro/eixo, janelas com RMS acima de 30\% da mediana. Essa é uma heurística de atividade, não uma confirmação do estado do motor. Das 2\,919 janelas completas somadas entre eixos, 101 foram excluídas; restam \textbf{2\,818}: 237 saudáveis por eixo e 703 danificadas em $x/y$, 701 em $z$. Sobras inferiores a 1~s não entram. Como a seleção é por eixo, os conjuntos não são inteiramente pareados.

Os tempos são crescentes e os dados finitos. Em F1/SV2 e F1/SV3 há um intervalo isolado de 2~ms na IMU; os resultados mantêm a FFT nominal de 1~kHz do Colab, uma aproximação a considerar. Não alinhamos IMU e piezo instante a instante. B1/B2 são mantidas e B3 fica limitada aos bins disponíveis até 500~Hz; B4 não é interpretada na IMU.

\begin{table}[!ht]
\caption{IMU após tratamento: médias por condição (S: saudável; D: danificada). RMS em $\U$, centroide em Hz. $d$ compara RMS entre condições no mesmo eixo.}
\label{tab:imu}\centering\small\setlength{\tabcolsep}{4pt}
\begin{tabular}{@{}lrrrrr@{}}\toprule
Eixo & RMS S & RMS D & Centr. S & Centr. D & $d$ RMS\\\midrule
$x$ & 2,267 & 2,536 & 280,7 & 290,2 & 0,58\\
$y$ & 1,964 & 2,666 & 272,9 & 242,0 & 1,21\\
$z$ & 5,502 & 7,295 & 319,8 & 346,3 & 1,42\\\bottomrule
\end{tabular}
\end{table}

\fig{07_notebook_imu_box}{IMU: distribuições após filtragem e seleção de atividade. Azul: saudável; vermelho: danificada. Os eixos são identificados por x, y e z em cada painel. As diferenças de intensidade e centroide dependem do eixo.}{fig:imu-box}

\fig{08_notebook_imu_espectros}{Espectros médios da IMU por eixo e condição, nos três painéis originais do notebook, com escalas comuns. Os pesos das componentes variam entre direções; a faixa observável termina em 500~Hz.}{fig:imu-espectros}

\subsection{Resposta à pergunta}
\textbf{Não, os eixos não têm as mesmas características.} Mesmo após atenuar a gravidade, o RMS saudável em $z$ é 2,4--2,8 vezes o de $x/y$ (Tabela~\ref{tab:imu}, Fig.~\ref{fig:imu-box}). A distribuição espectral também muda: $z$ dá mais peso à região de 390--400~Hz e tem centroide mais alto (Fig.~\ref{fig:imu-espectros}). A diferença entre condições não aparece igualmente: o aumento médio de RMS é cerca de 12\% em $x$, 36\% em $y$ e 33\% em $z$; o centroide cai em $y$ e sobe em $z$.

A potência em B3 aumenta nas três direções, com $d=1,48$, 0,97 e 1,44 para $x$, $y$ e $z$. O eixo $z$ se destaca no RMS e no centroide, enquanto $x$ tem o maior efeito em B3: não há um eixo universalmente melhor. A limitação de 500~Hz e o antialias de aquisição não verificado impedem descartar contaminação por frequências mais altas. Os efeitos dependem também da montagem e dos trechos retidos, sem validação de classificação.

\section{Conclusões}
A análise mostra dois aspectos complementares: a intensidade média muda, com sobreposição entre condições, e a potência se redistribui entre bandas, sobretudo B2 e B3. Os experimentos de duração e taxa mostram por que a leitura dos picos depende da observação e da amostragem. A comparação entre domínios mostra que janelas de mesma intensidade podem ter espectros distintos e que a posição dos picos acompanha a tensão da bateria, o que impede atribuí-la ao defeito. A IMU confirma que a direção altera intensidade e distribuição espectral. Esses achados justificam testar atributos temporais e bandas na próxima etapa, respeitando os poucos registros, a diversidade das falhas e a separação por aquisição.

O \href{https://colab.research.google.com/drive/1FdcHGuE1c6JbiSro7K0QYuGb-4w69KOJ?usp=sharing}{\textbf{Colab com código, gráficos e tabelas}} contém o percurso completo até a extração; não apresentamos resultados de classificação nesta entrega.

\begin{thebibliography}{1}
\bibitem{dataset} DronePropB, \emph{Ground Testing Dataset for Commercial Drones with Defective Propellers}, Mendeley Data, v.~1. DOI: \href{https://doi.org/10.17632/xkvfjmm8zg.1}{10.17632/xkvfjmm8zg.1}.
\end{thebibliography}
\end{document}
